% ---- score2 standard preamble (from: python3 score2_tools.py preamble) ----
\documentclass[10pt]{extarticle}
\usepackage[a4paper,margin=1.5cm]{geometry}
\usepackage[T1]{fontenc}
\usepackage{amsmath,amssymb}
\usepackage{graphicx}
\usepackage[export]{adjustbox}
\usepackage{multicol}
\usepackage{enumitem}
\usepackage{booktabs,array,tabularx}
\setlength{\columnsep}{0.7cm}
\setlength{\parindent}{0pt}
\setlength{\parskip}{3pt}
\pagestyle{empty}
\makeatletter
% figure registry (figdefs.tex is generated by the extract step)
\newcommand{\qfigdef}[4]{\@namedef{qf@file@#1}{#2}\@namedef{qf@w@#1}{#3}\@namedef{qf@wide@#1}{#4}}
\newcommand{\qf@check}[1]{\@ifundefined{qf@file@#1}{\PackageError{score2}{Unknown figure id #1}{}}{}}
\newcommand{\qf@img}[4][]{\includegraphics[width=#2\dimexpr\@nameuse{qf@w@#4}pt\relax,max width=#3,max height=0.45\textheight,#1]{\@nameuse{qf@file@#4}}}
% Full-width questions: finalize wraps every question that contains a wide figure in \qfullbegin...\qfullend,
% so the WHOLE question (stem, figures, options) is set across the page in reading order, then the two columns
% resume. The question's head - its number, stem and first figure (\qfullsplit marks the end of the head) -
% is kept on one page, so a number is never left at a page foot with its figure overleaf; the rest may break
% across pages normally, so no big gaps are left. Without a \qfullsplit, the whole question is kept together
% if it is under 45% of a page. Never write these yourself.
\newif\ifqf@full
\newif\ifqf@boxed
\newsavebox{\qf@box}
\newcommand{\qf@open}{\global\qf@fulltrue\par\global\qf@boxedtrue
  \setbox\qf@box\vbox\bgroup\hsize\textwidth\linewidth\textwidth}
\newcommand{\qf@emit}[1]{\par\egroup\global\qf@boxedfalse
  \ifdim\dimexpr\ht\qf@box+\dp\qf@box\relax<#1\textheight\noindent\usebox{\qf@box}\par\else\unvbox\qf@box\fi}
\newcommand{\qfullsplit}{\ifqf@boxed\qf@emit{0.85}\fi}
\newcommand{\qf@close}{\ifqf@boxed\qf@emit{0.45}\fi\global\qf@fullfalse}
\newcommand{\qfullbegin}{\end{multicols}\qf@open}
\newcommand{\qfullend}{\qf@close\begin{multicols}{2}}
\newcommand{\qfullstart}{\qf@open}
\newcommand{\qfulldone}{\qf@close}
% \qfig[scale]{id}: block figure, centred, directly where it is written (size set by the tool).
% Fallback if a wide figure is not inside a full-width question: break out of the columns in place - no floats.
\newcommand{\qfig}[2][1]{\qf@check{#2}%
  \ifqf@full\qf@blk{#1}{#2}\else
  \ifnum\@nameuse{qf@wide@#2}=1 \expandafter\@firstoftwo\else\expandafter\@secondoftwo\fi
  {\end{multicols}\begin{center}\qf@img{#1}{\textwidth}{#2}\end{center}\begin{multicols}{2}}%
  {\qf@blk{#1}{#2}}\fi}
\newcommand{\qf@blk}[2]{\par\smallskip{\centering\qf@img{#1}{\linewidth}{#2}\par}\smallskip}
% \qfigi[scale]{id}: inline figure (inside option lists / tables / minipages); never breaks columns.
\newcommand{\qfigi}[2][1]{\qf@check{#2}\qf@img[valign=c]{#1}{\linewidth}{#2}}
\makeatother
% structure macros - lint reads these to check completeness and figure placement
\newcommand{\papersec}[1]{\par\bigskip{\centering\Large\bfseries PAPER-#1\par}\medskip}
\newcommand{\partsec}[2]{\par\medskip\noindent{\large\bfseries PART-#1: #2}\par\smallskip}
\newcommand{\sect}[1]{\par\smallskip\noindent\textbf{#1}\par}
\newcommand{\qn}[1]{\par\smallskip\noindent\textbf{#1.}~}
\newcommand{\paraq}[2]{\par\smallskip\noindent\textbf{Paragraph for Questions #1 and #2}\par}
\newenvironment{qoptions}{\begin{itemize}[leftmargin=1.8em,itemsep=1pt,topsep=2pt,label={}]}{\end{itemize}}
\newcommand{\opt}[2]{\item[(#1)] #2}
% figure registry (generated by score2_tools.py finalize)
\qfigdef{p1_phy_s1_q01_1}{Score-I - Test 09 (16 Feb 2022)-img/p1_phy_s1_q01_1.png}{206.5}{0}
\qfigdef{p1_phy_s1_q02_1}{Score-I - Test 09 (16 Feb 2022)-img/p1_phy_s1_q02_1.png}{113.1}{0}
\qfigdef{p1_phy_s1_q03_1}{Score-I - Test 09 (16 Feb 2022)-img/p1_phy_s1_q03_1.png}{132.7}{0}
\qfigdef{p1_phy_s2_q05_1}{Score-I - Test 09 (16 Feb 2022)-img/p1_phy_s2_q05_1.png}{107.3}{0}
\qfigdef{p1_phy_s2_q06_1}{Score-I - Test 09 (16 Feb 2022)-img/p1_phy_s2_q06_1.png}{183.6}{0}
\qfigdef{p1_phy_s2_q09_1}{Score-I - Test 09 (16 Feb 2022)-img/p1_phy_s2_q09_1.png}{77.3}{0}
\qfigdef{p1_phy_s2_q10_1}{Score-I - Test 09 (16 Feb 2022)-img/p1_phy_s2_q10_1.png}{63.9}{0}
\qfigdef{p1_phy_s3_para3_4_1}{Score-I - Test 09 (16 Feb 2022)-img/p1_phy_s3_para3_4_1.png}{166.1}{0}
\qfigdef{p1_phy_s3_para5_6_1}{Score-I - Test 09 (16 Feb 2022)-img/p1_phy_s3_para5_6_1.png}{131.3}{0}
\qfigdef{p1_phy_s4_q08_1}{Score-I - Test 09 (16 Feb 2022)-img/p1_phy_s4_q08_1.png}{154.6}{0}
\qfigdef{p1_chem_s1_q01_1}{Score-I - Test 09 (16 Feb 2022)-img/p1_chem_s1_q01_1.png}{147.4}{0}
\qfigdef{p1_chem_s1_q01_A}{Score-I - Test 09 (16 Feb 2022)-img/p1_chem_s1_q01_A.png}{102.7}{0}
\qfigdef{p1_chem_s1_q01_B}{Score-I - Test 09 (16 Feb 2022)-img/p1_chem_s1_q01_B.png}{92.1}{0}
\qfigdef{p1_chem_s1_q01_C}{Score-I - Test 09 (16 Feb 2022)-img/p1_chem_s1_q01_C.png}{88.4}{0}
\qfigdef{p1_chem_s1_q01_D}{Score-I - Test 09 (16 Feb 2022)-img/p1_chem_s1_q01_D.png}{93.8}{0}
\qfigdef{p1_chem_s1_q02_A}{Score-I - Test 09 (16 Feb 2022)-img/p1_chem_s1_q02_A.png}{64.2}{0}
\qfigdef{p1_chem_s1_q02_B}{Score-I - Test 09 (16 Feb 2022)-img/p1_chem_s1_q02_B.png}{63.0}{0}
\qfigdef{p1_chem_s1_q02_C}{Score-I - Test 09 (16 Feb 2022)-img/p1_chem_s1_q02_C.png}{69.6}{0}
\qfigdef{p1_chem_s1_q02_D}{Score-I - Test 09 (16 Feb 2022)-img/p1_chem_s1_q02_D.png}{66.4}{0}
\qfigdef{p1_chem_s1_q03_1}{Score-I - Test 09 (16 Feb 2022)-img/p1_chem_s1_q03_1.png}{73.2}{0}
\qfigdef{p1_chem_s2_q05_1}{Score-I - Test 09 (16 Feb 2022)-img/p1_chem_s2_q05_1.png}{147.2}{0}
\qfigdef{p1_chem_s2_q05_A}{Score-I - Test 09 (16 Feb 2022)-img/p1_chem_s2_q05_A.png}{48.3}{0}
\qfigdef{p1_chem_s2_q05_B}{Score-I - Test 09 (16 Feb 2022)-img/p1_chem_s2_q05_B.png}{48.7}{0}
\qfigdef{p1_chem_s2_q05_C}{Score-I - Test 09 (16 Feb 2022)-img/p1_chem_s2_q05_C.png}{49.2}{0}
\qfigdef{p1_chem_s2_q05_D}{Score-I - Test 09 (16 Feb 2022)-img/p1_chem_s2_q05_D.png}{49.5}{0}
\qfigdef{p1_chem_s2_q06_A}{Score-I - Test 09 (16 Feb 2022)-img/p1_chem_s2_q06_A.png}{75.4}{0}
\qfigdef{p1_chem_s2_q06_B}{Score-I - Test 09 (16 Feb 2022)-img/p1_chem_s2_q06_B.png}{62.4}{0}
\qfigdef{p1_chem_s2_q06_C}{Score-I - Test 09 (16 Feb 2022)-img/p1_chem_s2_q06_C.png}{68.1}{0}
\qfigdef{p1_chem_s2_q06_D}{Score-I - Test 09 (16 Feb 2022)-img/p1_chem_s2_q06_D.png}{80.4}{0}
\qfigdef{p1_chem_s2_q08_A}{Score-I - Test 09 (16 Feb 2022)-img/p1_chem_s2_q08_A.png}{57.6}{0}
\qfigdef{p1_chem_s2_q08_B}{Score-I - Test 09 (16 Feb 2022)-img/p1_chem_s2_q08_B.png}{54.8}{0}
\qfigdef{p1_chem_s2_q08_C}{Score-I - Test 09 (16 Feb 2022)-img/p1_chem_s2_q08_C.png}{56.3}{0}
\qfigdef{p1_chem_s2_q08_D}{Score-I - Test 09 (16 Feb 2022)-img/p1_chem_s2_q08_D.png}{56.7}{0}
\qfigdef{p1_chem_s2_q10_1}{Score-I - Test 09 (16 Feb 2022)-img/p1_chem_s2_q10_1.png}{332.1}{1}
\qfigdef{p1_chem_s3_para1_2_1}{Score-I - Test 09 (16 Feb 2022)-img/p1_chem_s3_para1_2_1.png}{291.1}{1}
\qfigdef{p1_chem_s3_para3_4_1}{Score-I - Test 09 (16 Feb 2022)-img/p1_chem_s3_para3_4_1.png}{146.9}{0}
\qfigdef{p1_chem_s4_q08_1}{Score-I - Test 09 (16 Feb 2022)-img/p1_chem_s4_q08_1.png}{162.1}{0}
\qfigdef{p2_phy_s1_q01_1}{Score-I - Test 09 (16 Feb 2022)-img/p2_phy_s1_q01_1.png}{146.1}{0}
\qfigdef{p2_phy_s1_q03_1}{Score-I - Test 09 (16 Feb 2022)-img/p2_phy_s1_q03_1.png}{127.6}{0}
\qfigdef{p2_phy_s1_q06_1}{Score-I - Test 09 (16 Feb 2022)-img/p2_phy_s1_q06_1.png}{107.7}{0}
\qfigdef{p2_phy_s2_para7_8_1}{Score-I - Test 09 (16 Feb 2022)-img/p2_phy_s2_para7_8_1.png}{147.4}{0}
\qfigdef{p2_phy_s2_para9_10_1}{Score-I - Test 09 (16 Feb 2022)-img/p2_phy_s2_para9_10_1.png}{156.5}{0}
\qfigdef{p2_phy_s3_para1_2_1}{Score-I - Test 09 (16 Feb 2022)-img/p2_phy_s3_para1_2_1.png}{25.4}{0}
\qfigdef{p2_phy_s4_q07_1}{Score-I - Test 09 (16 Feb 2022)-img/p2_phy_s4_q07_1.png}{150.8}{0}
\qfigdef{p2_chem_s1_q01_A}{Score-I - Test 09 (16 Feb 2022)-img/p2_chem_s1_q01_A.png}{109.0}{0}
\qfigdef{p2_chem_s1_q01_B}{Score-I - Test 09 (16 Feb 2022)-img/p2_chem_s1_q01_B.png}{111.4}{0}
\qfigdef{p2_chem_s1_q01_C}{Score-I - Test 09 (16 Feb 2022)-img/p2_chem_s1_q01_C.png}{115.6}{0}
\qfigdef{p2_chem_s1_q01_D}{Score-I - Test 09 (16 Feb 2022)-img/p2_chem_s1_q01_D.png}{128.9}{0}
\qfigdef{p2_chem_s1_q02_1}{Score-I - Test 09 (16 Feb 2022)-img/p2_chem_s1_q02_1.png}{245.2}{0}
\qfigdef{p2_chem_s2_para7_8_1}{Score-I - Test 09 (16 Feb 2022)-img/p2_chem_s2_para7_8_1.png}{183.7}{0}
\qfigdef{p2_chem_s4_q07_1}{Score-I - Test 09 (16 Feb 2022)-img/p2_chem_s4_q07_1.png}{108.9}{0}
% ---- body skeleton ----
% \begin{document}
% \begin{center}{\Large\bfseries JEE (Main+Advanced) Score-II --- Test NN (DD Mon YYYY)}\end{center}
% \begin{multicols}{2}
%   \papersec{1}   (only for two-paper tests)
%   \partsec{1}{PHYSICS} \sect{Section-I (i): Single Correct} \qn{1} ... \qfig{p1_phy_q01_1} ...
% \end{multicols}
% \end{document}

\begin{document}
\raggedcolumns
\begin{center}{\Large\bfseries JEE (Main+Advanced) Score-I --- Test 09 (16 Feb 2022)}\end{center}
\begin{multicols}{2}
\papersec{1}
\raggedcolumns
\partsec{1}{PHYSICS}
\sect{SECTION-I(i)}
\qn{1} The smallest division on the main scale of a Vernier calipers is 0.1 cm. Ten divisions of the Vernier scale correspond to nine divisions of the main scale. The figure below on the left shows the reading of this calipers with no gap between its two jaws. The figure on the right shows the reading with a solid sphere held between the jaws. The correct diameter of the sphere is
\qfig{p1_phy_s1_q01_1}
\begin{qoptions}
\opt{A}{3.07 cm}
\opt{B}{3.11 cm}
\opt{C}{3.15 cm}
\opt{D}{3.17 cm}
\end{qoptions}

\qn{2} An ideal gas undergoes a four step cycle as shown in the $P$--$V$ diagram below. During this cycle, heat is absorbed by the gas in
\qfig{p1_phy_s1_q02_1}
\begin{qoptions}
\opt{A}{steps 1 and 2}
\opt{B}{steps 1 and 3}
\opt{C}{steps 1 and 4}
\opt{D}{steps 2 and 4}
\end{qoptions}

\qn{3} An extended object is placed at point O, 10 cm in front of a convex lens $L_1$ and a concave lens $L_2$ is placed 10 cm behind it, as shown in the figure. The radii of curvature of all the curved surfaces in both the lenses are 20 cm. The refractive index of both the lenses is 1.5. The total magnification of this lens system is
\qfig{p1_phy_s1_q03_1}
\begin{qoptions}
\opt{A}{0.4}
\opt{B}{0.8}
\opt{C}{1.3}
\opt{D}{1.6}
\end{qoptions}

\qn{4} A heavy nucleus Q of half-life 20 minutes undergoes alpha-decay with probability of 60\% and beta-decay with probability of 40\%. Initially, the number of Q nuclei is 1000. The number of alpha-decays of Q in the first one hour is
\begin{qoptions}
\opt{A}{50}
\opt{B}{75}
\opt{C}{350}
\opt{D}{525}
\end{qoptions}

\sect{SECTION-I(ii)}
\qn{5} A horizontal force $F$ is applied at the center of mass of a cylindrical object of mass $m$ and radius $R$, perpendicular to its axis as shown in the figure. The coefficient of friction between the object and the ground is $\mu$. The center of mass of the object has an acceleration $a$. The acceleration due to gravity is $g$. Given that the object rolls without slipping, which of the following statement(s) is(are) correct?
\qfig{p1_phy_s2_q05_1}
\begin{qoptions}
\opt{A}{For the same $F$, the value of $a$ does not depend on whether the cylinder is solid or hollow}
\opt{B}{For a solid cylinder, the maximum possible value of $a$ is $2\mu g$}
\opt{C}{The magnitude of the frictional force on the object due to the ground is always $\mu mg$}
\opt{D}{For a thin-walled hollow cylinder, $a=\dfrac{F}{2m}$}
\end{qoptions}

\qn{6} A wide slab consisting of two media of refractive indices $n_1$ and $n_2$ is placed in air as shown in the figure. A ray of light is incident from medium $n_1$ to $n_2$ at an angle $\theta$, where $\sin\theta$ is slightly larger than $1/n_1$. Take refractive index of air as 1. Which of the following statement(s) is(are) correct?
\qfig{p1_phy_s2_q06_1}
\begin{qoptions}
\opt{A}{The light ray enters air if $n_2=n_1$}
\opt{B}{The light ray is finally reflected back into the medium of refractive index $n_1$ if $n_2<n_1$}
\opt{C}{The light ray is finally reflected back into the medium of refractive index $n_1$ if $n_2>n_1$}
\opt{D}{The light ray is reflected back into the medium of refractive index $n_1$ if $n_2=1$}
\end{qoptions}

\qn{7} A particle of mass $M=0.2$ kg is initially at rest in the $xy$-plane at a point ($x=-\ell$, $y=-h$), where $\ell=10$ m and $h=1$ m. The particle is accelerated at time $t=0$ with a constant acceleration $a=10\ \text{m/s}^2$ along the positive $x$-direction. Its angular momentum and torque with respect to the origin, in SI units, are represented by $\vec L$ and $\vec\tau$, respectively. $\hat i$, $\hat j$ and $\hat k$ are unit vectors along the positive $x$, $y$ and $z$-directions, respectively. If $\hat k=\hat i\times\hat j$ then which of the following statement(s) is(are) correct ?
\begin{qoptions}
\opt{A}{The particle arrives at the point ($x=\ell$, $y=-h$) at time $t=2$ s.}
\opt{B}{$\vec\tau=2\hat k$ when the particle passes through the point ($x=\ell$, $y=-h$)}
\opt{C}{$\vec L=4\hat k$ when the particle passes through the point ($x=\ell$, $y=-h$)}
\opt{D}{$\vec\tau=\hat k$ when the particle passes through the point ($x=0$, $y=-h$)}
\end{qoptions}

\qn{8} Which of the following statement(s) is(are) correct about the spectrum of hydrogen atom ?
\begin{qoptions}
\opt{A}{The ratio of the longest wavelength to the shortest wavelength in Balmer series is $9/5$}
\opt{B}{There is an overlap between the wavelength ranges of Balmer and Paschen series.}
\opt{C}{The wavelengths of Lyman series are given by $\left(1+\dfrac{1}{m^2}\right)\lambda_0$, where $\lambda_0$ is the shortest wavelength of Lyman series and $m$ is an integer}
\opt{D}{The wavelength ranges of Lyman and Balmer series do not overlap}
\end{qoptions}

\qn{9} A long straight wire carries a current, $I=2$ ampere. A semi-circular conducting rod is placed beside it on two conducting parallel rails of negligible resistance. Both the rails are parallel to the wire. The wire, the rod and the rails lie in the same horizontal plane, as shown in the figure. Two ends of the semi-circular rod are at distances 1 cm and 4 cm from the wire. At time $t=0$, the rod starts moving on the rails with a speed $v=3.0$ m/s (see the figure). A resistor $R=1.4\ \Omega$ and a capacitor $C_0=5.0\ \mu$F are connected in series between the rails. At time $t=0$, $C_0$ is uncharged. Which of the following statement(s) is(are) correct ? [$\mu_0=4\pi\times10^{-7}$ SI units. Take $\ell n\,2=0.7$]
\qfig{p1_phy_s2_q09_1}
\begin{qoptions}
\opt{A}{Maximum current through $R$ is $1.2\times10^{-6}$ ampere}
\opt{B}{Maximum current through $R$ is $3.8\times10^{-6}$ ampere}
\opt{C}{Maximum charge on capacitor $C_0$ is $8.4\times10^{-12}$ coulomb}
\opt{D}{Maximum charge on capacitor $C_0$ is $2.4\times10^{-12}$ coulomb}
\end{qoptions}

\qn{10} A cylindrical tube, with its base as shown in the figure, is filled with water. It is moving down with a constant acceleration $a$ along a fixed inclined plane with angle $\theta=45^\circ$. $P_1$ and $P_2$ are pressures at points \textbf{1} and \textbf{2}, respectively, located at the base of the tube. Let $\beta=(P_1-P_2)/(\rho g d)$, where $\rho$ is density of water, $d$ is the inner diameter of the tube and $g$ is the acceleration due to gravity. Which of the following statement(s) is(are) correct ?
\qfig{p1_phy_s2_q10_1}
\begin{qoptions}
\opt{A}{$\beta=0$ when $a=g/\sqrt2$}
\opt{B}{$\beta>0$ when $a=g/\sqrt2$}
\opt{C}{$\beta=\dfrac{\sqrt2-1}{\sqrt2}$ when $a=g/2$}
\opt{D}{$\beta=\dfrac{1}{\sqrt2}$ when $a=g/2$}
\end{qoptions}

\sect{SECTION-II(i)}
\paraq{1}{2}
A projectile is thrown from a point O on the ground at an angle $45^\circ$ from the vertical and with a speed $5\sqrt2$ m/s. The projectile at the highest point of its trajectory splits into two equal parts. One part falls vertically down to the ground, 0.5 s after the splitting. The other part, $t$ seconds after the splitting, falls to the ground at a distance $x$ meters from the point O. The acceleration due to gravity $g=10\ \text{m/s}^2$.

\qn{1} The value of $t$ is \underline{\hspace{2.5em}}.

\qn{2} The value of $x$ is \underline{\hspace{2.5em}}.

\paraq{3}{4}
In the circuit shown below, the switch S is connected to position P for a long time so that the charge on the capacitor becomes $q_1\ \mu$C. Then S is switched to position Q. After a long time, the charge on the capacitor is $q_2\ \mu$C.
\qfig{p1_phy_s3_para3_4_1}

\qn{3} The magnitude of $q_1$ is \underline{\hspace{2.5em}}.

\qn{4} The magnitude of $q_2$ is \underline{\hspace{2.5em}}.

\paraq{5}{6}
Two point charges $-Q$ and $+Q/\sqrt3$ are placed in the $xy$-plane at the origin $(0,0)$ and a point $(2,0)$, respectively, as shown in the figure. This results in an equipotential circle of radius $R$ and potential $V=0$ in the $xy$-plane with its center at $(b,0)$. All lengths are measured in meters.
\qfig{p1_phy_s3_para5_6_1}

\qn{5} The value of $R$ is \underline{\hspace{2.5em}} meter.

\qn{6} The value of $b$ is \underline{\hspace{2.5em}} meter.

\sect{SECTION-II(ii)}
\qn{7} An $\alpha$-particle (mass 4 amu) and a singly charged sulfur ion (mass 32 amu) are initially at rest. They are accelerated through a potential $V$ and then allowed to pass into a region of uniform magnetic field which is normal to the velocities of the particles. Within this region, the $\alpha$-particle and the sulfur ion move in circular orbits of radii $r_\alpha$ and $r_s$, respectively. The ratio $(r_s/r_\alpha)$ is \underline{\hspace{2.5em}}.

\qn{8} A thin rod of mass $M$ and length $a$ is free to rotate in horizontal plane about a fixed vertical axis passing through point O. A thin circular disc of mass $M$ and of radius $a/4$ is pivoted on this rod with its center at a distance $a/4$ from the free end so that it can rotate freely about its vertical axis, as shown in the figure. Assume that both the rod and the disc have uniform density and they remain horizontal during the motion. An outside stationary observer finds the rod rotating with an angular velocity $\Omega$ and the disc rotating about its vertical axis with angular velocity $4\Omega$. The total angular momentum of the system about the point O is $\left(\dfrac{Ma^2\Omega}{48}\right)n$. The value of $n$ is \underline{\hspace{2.5em}}.
\qfig{p1_phy_s4_q08_1}

\qn{9} A small object is placed at the center of a large evacuated hollow spherical container. Assume that the container is maintained at 0 K. At time $t=0$, the temperature of the object is 200 K. The temperature of the object becomes 100 K at $t=t_1$ and 50 K at $t=t_2$. Assume the object and the container to be ideal black bodies. The heat capacity of the object does not depend on temperature. The ratio $(t_2/t_1)$ is \underline{\hspace{2.5em}}.

\raggedcolumns
\partsec{2}{CHEMISTRY}
\sect{SECTION-I(i)}
\qn{1} The major product formed in the following reaction is
\qfig{p1_chem_s1_q01_1}
\begin{qoptions}
\opt{A}{\qfigi{p1_chem_s1_q01_A}}
\opt{B}{\qfigi{p1_chem_s1_q01_B}}
\opt{C}{\qfigi{p1_chem_s1_q01_C}}
\opt{D}{\qfigi{p1_chem_s1_q01_D}}
\end{qoptions}

\qn{2} Among the following, the conformation that corresponds to the most stable conformation of \textit{meso}--butane--2,3--diol is --
\begin{qoptions}
\opt{A}{\qfigi{p1_chem_s1_q02_A}}
\opt{B}{\qfigi{p1_chem_s1_q02_B}}
\opt{C}{\qfigi{p1_chem_s1_q02_C}}
\opt{D}{\qfigi{p1_chem_s1_q02_D}}
\end{qoptions}

\qn{3} For the given close packed structure of a salt made of cation \textbf{X} and anion \textbf{Y} shown below (ions of only one face are shown for clarity), the packing fraction is approximately (packing fraction $=\dfrac{\text{Packing efficiency}}{100}$)
\qfig{p1_chem_s1_q03_1}
\begin{qoptions}
\opt{A}{0.74}
\opt{B}{0.63}
\opt{C}{0.52}
\opt{D}{0.48}
\end{qoptions}

\qn{4} The calculated spin only magnetic moments of $[\mathrm{Cr(NH_3)_6}]^{3+}$ and $[\mathrm{CuF_6}]^{3-}$ in BM, respectively, are (Atomic numbers of Cr and Cu are 24 and 29, respectively)
\begin{qoptions}
\opt{A}{3.87 and 2.84}
\opt{B}{4.90 and 1.73}
\opt{C}{3.87 and 1.73}
\opt{D}{4.90 and 2.84}
\end{qoptions}

\sect{SECTION-I(ii)}
\qn{5} Given
\qfig{p1_chem_s2_q05_1}
The compound(s), which on reaction with $\mathrm{HNO_3}$ will give the product having degree of rotation, $[\alpha]_D=-52.7^\circ$ is (are)
\begin{qoptions}
\opt{A}{\qfigi{p1_chem_s2_q05_A}}
\opt{B}{\qfigi{p1_chem_s2_q05_B}}
\opt{C}{\qfigi{p1_chem_s2_q05_C}}
\opt{D}{\qfigi{p1_chem_s2_q05_D}}
\end{qoptions}

\qn{6} The reaction of \textbf{Q} with PhSNa yields an organic compound (major product) that gives positive Carius test on treatment with $\mathrm{Na_2O_2}$ followed by addition of $\mathrm{BaCl_2}$. The correct option(s) for \textbf{Q} is (are).
\begin{qoptions}
\opt{A}{\qfigi{p1_chem_s2_q06_A}}
\opt{B}{\qfigi{p1_chem_s2_q06_B}}
\opt{C}{\qfigi{p1_chem_s2_q06_C}}
\opt{D}{\qfigi{p1_chem_s2_q06_D}}
\end{qoptions}

\qn{7} The correct statement(s) related to colloids is(are)
\begin{qoptions}
\opt{A}{The process of precipitating colloidal sol by an electrolyte is called peptization.}
\opt{B}{Colloidal solution freezes at higher temperature than the true solution at the same concentration.}
\opt{C}{Surfactants form micelle above critical micelle concentration (CMC). CMC depends on temperature}
\opt{D}{Micelles are macromolecular colloids.}
\end{qoptions}

\qn{8} An ideal gas undergoes a reversible isothermal expansion from state \textbf{I} to state \textbf{II} followed by a reversible adiabatic expansion from state \textbf{II} to state \textbf{III}. The correct plot(s) representing the changes from state \textbf{I} to state \textbf{III} is(are)\\
($p$ : pressure, $V$ : volume, $T$ : temperature, $H$ : enthalpy, $S$ : entropy)
\begin{qoptions}
\opt{A}{\qfigi{p1_chem_s2_q08_A}}
\opt{B}{\qfigi{p1_chem_s2_q08_B}}
\opt{C}{\qfigi{p1_chem_s2_q08_C}}
\opt{D}{\qfigi{p1_chem_s2_q08_D}}
\end{qoptions}

\qn{9} The correct statement(s) related to the metal extraction processes is(are)
\begin{qoptions}
\opt{A}{A mixture of PbS and PbO undergoes self-reduction to produce Pb and $\mathrm{SO_2}$.}
\opt{B}{In the extraction process of copper from copper pyrites, silica is added to produce copper silicate}
\opt{C}{Partial oxidation of sulphide ore of copper by roasting, followed by self-reduction produces blister copper.}
\opt{D}{In cyanide process, zinc powder is utilized to precipitate gold from $\mathrm{Na[Au(CN)_2]}$}
\end{qoptions}

\qfullbegin
\qn{10} A mixture of two salts is used to prepare a solution S, which gives the following results :
\qfig{p1_chem_s2_q10_1}
\qfullsplit
The correct option(s) for the salt mixture is(are)
\begin{qoptions}
\opt{A}{$\mathrm{Pb(NO_3)_2}$ and $\mathrm{Zn(NO_3)_2}$}
\opt{B}{$\mathrm{Pb(NO_3)_2}$ and $\mathrm{Bi(NO_3)_3}$}
\opt{C}{$\mathrm{AgNO_3}$ and $\mathrm{Bi(NO_3)_3}$}
\opt{D}{$\mathrm{Pb(NO_3)_2}$ and $\mathrm{Hg(NO_3)_2}$}
\end{qoptions}
\sect{SECTION-II(i)}
\paraq{1}{2}
For the following reaction scheme, percentage yields are given along the arrow :
\qfig{p1_chem_s3_para1_2_1}
\qfullsplit
$\mathbf{x}$ g and $\mathbf{y}$ g are mass of \textbf{R} and \textbf{U}, respectively.\\
(Use : Molar mass (in g mol$^{-1}$) of H, C and O as 1, 12 and 16, respectively)
\qfullend

\qn{1} The value of $\mathbf{x}$ is \underline{\hspace{2.5em}}.

\qn{2} The value of $\mathbf{y}$ is \underline{\hspace{2.5em}}.

\paraq{3}{4}
For the reaction $\mathrm{X(s)}\rightleftharpoons \mathrm{Y(s)}+\mathrm{Z(g)}$, the plot of $\ln\dfrac{p_z}{p^{\ominus}}$ versus $\dfrac{10^4}{T}$ is given below (in solid line), where $p_z$ is the pressure (in bar) of the gas \textbf{Z} at temperature $T$ and $p^{\ominus}=1$ bar.
\qfig{p1_chem_s3_para3_4_1}
(Given, $\dfrac{d(\ln K)}{d\left(\frac1T\right)}=-\dfrac{\Delta H^{\ominus}}{R}$, where the equilibrium constant, $K=\dfrac{p_z}{p^{\ominus}}$ and the gas constant, $R=8.314\ \mathrm{J\,K^{-1}\,mol^{-1}}$)

\qn{3} The value of standard enthalpy, $\Delta H^{\ominus}$ (in kJ mol$^{-1}$) for the reaction is \underline{\hspace{2.5em}}.

\qn{4} The value of $\Delta S^{\ominus}$ (in J K$^{-1}$ mol$^{-1}$) for the given reaction, at 1000 K is \underline{\hspace{2.5em}}.

\paraq{5}{6}
The boiling point of water in a 0.1 molal silver nitrate solution (solution \textbf{A}) is $\mathbf{x}\ ^\circ$C. To this solution \textbf{A}, an equal volume of 0.1 molal aqueous barium chloride solution is added to make a new solution \textbf{B}. The difference in the boiling points of water in the two solutions \textbf{A} and \textbf{B} is $\mathbf{y}\times10^{-2}\ ^\circ$C.\\
(Assume : Densities of the solutions \textbf{A} and \textbf{B} are the same as that of water and the soluble salts dissociate completely.)\\
Use: Molal elevation constant (Ebullioscopic Constant), $K_b=0.5\ \mathrm{K\,kg\,mol^{-1}}$; Boiling point of pure water as $100^\circ$C.

\qn{5} The value of $\mathbf{x}$ is \underline{\hspace{2.5em}}.

\qn{6} The value of $|\mathbf{y}|$ is \underline{\hspace{2.5em}}.

\sect{SECTION-II(ii)}
\qn{7} The maximum number of possible isomers (including stereoisomers) which may be formed on \textit{mono}-bromination of 1-methylcyclohex-1-ene using $\mathrm{Br_2}$ and UV light is \underline{\hspace{2.5em}}

\qn{8} In the reaction given below, the total number of atoms having $sp^2$ hybridization in the major product P is \underline{\hspace{2.5em}}
\qfig{p1_chem_s4_q08_1}

\qn{9} The total number of possible isomers for $[\mathrm{Pt(NH_3)_4Cl_2}]\mathrm{Br_2}$ is \underline{\hspace{2.5em}}

\raggedcolumns
\partsec{3}{MATHEMATICS}
\sect{SECTION-I(i)}
\qn{1} Consider a triangle $\Delta$ whose two sides lie on the $x$-axis and the line $x+y+1=0$. If the orthocenter of $\Delta$ is $(1,1)$, then the equation of the circle passing through the vertices of the triangle $\Delta$ is
\begin{qoptions}
\opt{A}{$x^2+y^2-3x+y=0$}
\opt{B}{$x^2+y^2+x+3y=0$}
\opt{C}{$x^2+y^2+2y-1=0$}
\opt{D}{$x^2+y^2+x+y=0$}
\end{qoptions}

\qn{2} The area of the region\\
{\small$\left\{(x,y):\ 0\le x\le\dfrac94,\ 0\le y\le1,\ x\ge3y,\ x+y\ge2\right\}$}\\ is
\begin{qoptions}
\opt{A}{$\dfrac{11}{32}$}
\opt{B}{$\dfrac{35}{96}$}
\opt{C}{$\dfrac{37}{96}$}
\opt{D}{$\dfrac{13}{32}$}
\end{qoptions}

\qn{3} Consider three sets $E_1=\{1,2,3\}$, $F_1=\{1,3,4\}$ and $G_1=\{2,3,4,5\}$. Two elements are chosen at random, without replacement, from the set $E_1$, and let $S_1$ denote the set of these chosen elements.\\
Let $E_2=E_1-S_1$ and $F_2=F_1\cup S_1$. Now two elements are chosen at random, without replacement, from the set $F_2$ and let $S_2$ denote the set of these chosen elements.\\
Let $G_2=G_1\cup S_2$. Finally, two elements are chosen at random, without replacement, from the set $G_2$ and let $S_3$ denote the set of these chosen elements.\\
Let $E_3=E_2\cup S_3$. Given that $E_1=E_3$, let $p$ be the conditional probability of the event $S_1=\{1,2\}$. Then the value of $p$ is
\begin{qoptions}
\opt{A}{$\dfrac15$}
\opt{B}{$\dfrac35$}
\opt{C}{$\dfrac12$}
\opt{D}{$\dfrac25$}
\end{qoptions}

\qn{4} Let $\theta_1,\theta_2,\dots,\theta_{10}$ be positive valued angles (in radian) such that $\theta_1+\theta_2+\dots+\theta_{10}=2\pi$. Define the complex numbers $z_1=e^{i\theta_1}$, $z_k=z_{k-1}e^{i\theta_k}$ for $k=2,3,\dots,10$, where $i=\sqrt{-1}$. Consider the statements P and Q given below :\\
P : $|z_2-z_1|+|z_3-z_2|+\dots+|z_{10}-z_9|+|z_1-z_{10}|\le2\pi$\\
Q : $|z_2^2-z_1^2|+|z_3^2-z_2^2|+\dots+|z_{10}^2-z_9^2|+|z_1^2-z_{10}^2|\le4\pi$\\
Then,
\begin{qoptions}
\opt{A}{P is \textbf{TRUE} and Q is \textbf{FALSE}}
\opt{B}{Q is \textbf{TRUE} and P is \textbf{FALSE}}
\opt{C}{both P and Q are \textbf{TRUE}}
\opt{D}{both P and Q are \textbf{FALSE}}
\end{qoptions}

\sect{SECTION-I(ii)}
\qn{5} For any $3\times3$ matrix $M$, let $|M|$ denote the determinant of $M$. Let
\begin{gather*}
E=\begin{bmatrix}1&2&3\\2&3&4\\8&13&18\end{bmatrix},\quad
P=\begin{bmatrix}1&0&0\\0&0&1\\0&1&0\end{bmatrix}\\
\text{and}\quad F=\begin{bmatrix}1&3&2\\8&18&13\\2&4&3\end{bmatrix}.
\end{gather*}
If $Q$ is a nonsingular matrix of order $3\times3$, then which of the following statements is (are) \textbf{TRUE} ?
\begin{qoptions}
\opt{A}{$F=PEP$ and $P^2=\begin{bmatrix}1&0&0\\0&1&0\\0&0&1\end{bmatrix}$}
\opt{B}{$|EQ+PFQ^{-1}|=|EQ|+|PFQ^{-1}|$}
\opt{C}{$\left|(EF)^3\right|>|EF|^2$}
\opt{D}{Sum of the diagonal entries of $P^{-1}EP+F$ is equal to the sum of diagonal entries of $E+P^{-1}FP$}
\end{qoptions}

\qn{6} Let $f:\mathbb R\to\mathbb R$ be defined by $f(x)=\dfrac{x^2-3x-6}{x^2+2x+4}$. Then which of the following statements is (are) \textbf{TRUE} ?
\begin{qoptions}
\opt{A}{$f$ is decreasing in the interval $(-2,-1)$}
\opt{B}{$f$ is increasing in the interval $(1,2)$}
\opt{C}{$f$ is onto}
\opt{D}{Range of $f$ is $\left[-\dfrac32,2\right]$}
\end{qoptions}

\qn{7} Let $E$, $F$ and $G$ be three events having probabilities $P(E)=\dfrac18$, $P(F)=\dfrac16$ and $P(G)=\dfrac14$, and let $P(E\cap F\cap G)=\dfrac1{10}$. For any event $H$, if $H^C$ denotes its complement, then which of the following statements is(are) \textbf{TRUE} ?
\begin{qoptions}
\opt{A}{$P(E\cap F\cap G^C)\le\dfrac1{40}$}
\opt{B}{$P(E^C\cap F\cap G)\le\dfrac1{15}$}
\opt{C}{$P(E\cup F\cup G)\le\dfrac{13}{24}$}
\opt{D}{$P(E^C\cap F^C\cap G^C)\le\dfrac5{12}$}
\end{qoptions}

\qn{8} For any $3\times3$ matrix $M$, let $|M|$ denote the determinant of $M$. Let $I$ be the $3\times3$ identity matrix. Let $E$ and $F$ be two $3\times3$ matrices such that $(I-EF)$ is invertible. If $G=(I-EF)^{-1}$, then which of the following statements is (are) \textbf{TRUE} ?
\begin{qoptions}
\opt{A}{$|FE|=|I-FE|\,|FGE|$}
\opt{B}{$|I-FE|\,(I+FGE)=I$}
\opt{C}{$EFG=GEF$}
\opt{D}{$(I-FE)(I-FGE)=I$}
\end{qoptions}

\qn{9} For any positive integer $n$, let $S_n:(0,\infty)\to\mathbb R$ be defined by $S_n(x)=\displaystyle\sum_{k=1}^{n}\cot^{-1}\!\left(\dfrac{1+k(k+1)x^2}{x}\right)$, where for any $x\in\mathbb R$, $\cot^{-1}x\in(0,\pi)$ and $\tan^{-1}(x)\in\left(-\dfrac\pi2,\dfrac\pi2\right)$. Then which of the following statements is (are) \textbf{TRUE} ?
\begin{qoptions}
\opt{A}{$S_{10}(x)=\dfrac\pi2-\tan^{-1}\!\left(\dfrac{1+11x^2}{10x}\right)$, for all $x>0$}
\opt{B}{$\displaystyle\lim_{n\to\infty}\cot(S_n(x))=x$, for all $x>0$}
\opt{C}{The equation $S_3(x)=\dfrac\pi4$ has a root in $(0,\infty)$}
\opt{D}{$\tan(S_n(x))\le\dfrac12$, for all $n\ge1$ and $x>0$}
\end{qoptions}

\qn{10} For any complex number $w=c+id$, let $\arg(w)\in(-\pi,\pi]$, where $i=\sqrt{-1}$. Let $\alpha$ and $\beta$ be real numbers such that for all complex numbers $z=x+iy$ satisfying $\arg\!\left(\dfrac{z+\alpha}{z+\beta}\right)=\dfrac\pi4$, the ordered pair $(x,y)$ lies on the circle $x^2+y^2+5x-3y+4=0$. Then which of the following statements is (are) \textbf{TRUE} ?
\begin{qoptions}
\opt{A}{$\alpha=-1$}
\opt{B}{$\alpha\beta=4$}
\opt{C}{$\alpha\beta=-4$}
\opt{D}{$\beta=4$}
\end{qoptions}

\sect{SECTION-II(i)}
\paraq{1}{2}
Three numbers are chosen at random, one after another with replacement, from the set $S=\{1,2,3,\dots,100\}$. Let $p_1$ be the probability that the maximum of chosen numbers is at least 81 and $p_2$ be the probability that the minimum of chosen numbers is at most 40.

\qn{1} The value of $\dfrac{625}{4}p_1$ is \underline{\hspace{3em}}.

\qn{2} The value of $\dfrac{125}{4}p_2$ is \underline{\hspace{3em}}.

\paraq{3}{4}
Let $\alpha$, $\beta$ and $\gamma$ be real numbers such that the system of linear equations
\begin{align*}
x+2y+3z&=\alpha\\
4x+5y+6z&=\beta\\
7x+8y+9z&=\gamma-1
\end{align*}
is consistent. Let $|M|$ represent the determinant of the matrix $M=\begin{bmatrix}\alpha&2&\gamma\\\beta&1&0\\-1&0&1\end{bmatrix}$.
Let $P$ be the plane containing all those $(\alpha,\beta,\gamma)$ for which the above system of linear equations is consistent, and $D$ be the \textbf{square} of the distance of the point $(0,1,0)$ from the plane $P$.

\qn{3} The value of $|M|$ is \underline{\hspace{3em}}.

\qn{4} The value of $D$ is \underline{\hspace{3em}}.

\paraq{5}{6}
Consider the lines $L_1$ and $L_2$ defined by $L_1:x\sqrt2+y-1=0$ and $L_2:x\sqrt2-y+1=0$.\\
For a fixed constant $\lambda$, let $C$ be the locus of a point $P$ such that the product of the distance of $P$ from $L_1$ and the distance of $P$ from $L_2$ is $\lambda^2$. The line $y=2x+1$ meets $C$ at two points $R$ and $S$, where the distance between $R$ and $S$ is $\sqrt{270}$.\\
Let the perpendicular bisector of $RS$ meet $C$ at two distinct points $R'$ and $S'$. Let $D$ be the \textbf{square} of the distance between $R'$ and $S'$.

\qn{5} The value of $\lambda^2$ is \underline{\hspace{3em}}.

\qn{6} The value of $D$ is \underline{\hspace{3em}}.

\sect{SECTION-II(ii)}
\qn{7} For $x\in\mathbb R$, then number of real roots of the equation $3x^2-4|x^2-1|+x-1=0$ is \underline{\hspace{3em}}.

\qn{8} In a triangle $ABC$, let $AB=\sqrt{23}$, $BC=3$ and $CA=4$. Then the value of $\dfrac{\cot A+\cot C}{\cot B}$ is \underline{\hspace{3em}}.

\qn{9} Let $\vec u$, $\vec v$ and $\vec w$ be vectors in three-dimensional space, where $\vec u$ and $\vec v$ are unit vectors which are not perpendicular to each other and $\vec u\cdot\vec w=1$, $\vec v\cdot\vec w=1$, $\vec w\cdot\vec w=4$. If the volume of the parallelopiped, whose adjacent sides are represented by the vectors $\vec u$, $\vec v$ and $\vec w$, is $\sqrt2$, then the value of $|3\vec u+5\vec v|$ is \underline{\hspace{3em}}.

\raggedcolumns
\papersec{2}
\partsec{1}{PHYSICS}
\sect{SECTION-I(i)}
\qn{1} One end of a horizontal uniform beam of weight $W$ and length $L$ is hinged on a vertical wall at point $O$ and its other end is supported by a light inextensible rope. The other end of the rope is fixed at point $Q$, at a height $L$ above the hinge at point $O$. A block of weight $\alpha W$ is attached at the point $P$ of the beam, as shown in the figure (not to scale). The rope can sustain a maximum tension of $(2\sqrt2)W$. Which of the following statement(s) is(are) correct ?
\qfig{p2_phy_s1_q01_1}
\begin{qoptions}
\opt{A}{The vertical component of reaction force at $O$ does \textbf{not} depend on $\alpha$}
\opt{B}{The horizontal component of reaction force at $O$ is equal to $W$ for $\alpha=0.5$}
\opt{C}{The tension in the rope is $2W$ for $\alpha=0.5$}
\opt{D}{The rope breaks if $\alpha>1.5$}
\end{qoptions}

\qn{2} A source, approaching with speed $u$ towards the open end of a stationary pipe of length $L$, is emitting a sound of frequency $f_s$. The farther end of the pipe is closed. The speed of sound in air is $v$ and $f_0$ is the fundamental frequency of the pipe. For which of the following combination(s) of $u$ and $f_s$, will the sound reaching the pipe lead to a resonance ?
\begin{qoptions}
\opt{A}{$u=0.8\,v$ and $f_s=f_0$}
\opt{B}{$u=0.8\,v$ and $f_s=2f_0$}
\opt{C}{$u=0.8\,v$ and $f_s=0.5\,f_0$}
\opt{D}{$u=0.5\,v$ and $f_s=1.5\,f_0$}
\end{qoptions}

\qn{3} For a prism of prism angle $\theta=60^\circ$, the refractive indices of the left half and the right half are, respectively, $n_1$ and $n_2$ ($n_2\ge n_1$) as shown in the figure. The angle of incidence $i$ is chosen such that the incident light rays will have minimum deviation if $n_1=n_2=n=1.5$. For the case of unequal refractive indices, $n_1=n$ and $n_2=n+\Delta n$ (where $\Delta n\ll n$), the angle of emergence $e=i+\Delta e$. Which of the following statement(s) is (are) correct ?
\qfig{p2_phy_s1_q03_1}
\begin{qoptions}
\opt{A}{The value of $\Delta e$ (in radians) is greater than that of $\Delta n$}
\opt{B}{$\Delta e$ is proportional to $\Delta n$}
\opt{C}{$\Delta e$ lies between 2.0 and 3.0 milliradians, if $\Delta n=2.8\times10^{-3}$}
\opt{D}{$\Delta e$ lies between 1.0 and 1.6 milliradians, if $\Delta n=2.8\times10^{-3}$}
\end{qoptions}

\qn{4} A physical quantity $\vec S$ is defined as $\vec S=(\vec E\times\vec B)/\mu_0$, where $\vec E$ is electric field, $\vec B$ is magnetic field and $\mu_0$ is the permeability of free space. The dimensions of $\vec S$ are the same as the dimensions of which of the following quantity(ies) ?
\begin{qoptions}
\opt{A}{$\dfrac{\text{Energy}}{\text{charge}\times\text{current}}$}
\opt{B}{$\dfrac{\text{Force}}{\text{Length}\times\text{Time}}$}
\opt{C}{$\dfrac{\text{Energy}}{\text{Volume}}$}
\opt{D}{$\dfrac{\text{Power}}{\text{Area}}$}
\end{qoptions}

\qn{5} A heavy nucleus $N$, at rest, undergoes fission $N\to P+Q$, where $P$ and $Q$ are two lighter nuclei. Let $\delta=M_N-M_P-M_Q$, where $M_P$, $M_Q$ and $M_N$ are the masses of $P$, $Q$ and $N$, respectively. $E_P$ and $E_Q$ are the kinetic energies of $P$ and $Q$, respectively. The speed of $P$ and $Q$ are $v_P$ and $v_Q$, respectively. If $c$ is the speed of light, which of the following statement(s) is(are) correct ?
\begin{qoptions}
\opt{A}{$E_P+E_Q=c^2\delta$}
\opt{B}{$E_P=\left(\dfrac{M_P}{M_P+M_Q}\right)c^2\delta$}
\opt{C}{$\dfrac{v_P}{v_Q}=\dfrac{M_Q}{M_P}$}
\opt{D}{The magnitude of momentum for $P$ as well as $Q$ is $c\sqrt{2\mu\delta}$, where $\mu=\dfrac{M_PM_Q}{(M_P+M_Q)}$}
\end{qoptions}

\qn{6} Two concentric circular loops, one of radius $R$ and the other of radius $2R$, lie in the $xy$-plane with the origin as their common center, as shown in the figure. The smaller loop carries current $I_1$ in the anti-clockwise direction and the larger loop carries current $I_2$ in the clockwise direction, with $I_2>2I_1$. $\vec B(x,y)$ denotes the magnetic field at a point $(x,y)$ in the $xy$-plane. Which of the following statement(s) is(are) correct?
\qfig{p2_phy_s1_q06_1}
\begin{qoptions}
\opt{A}{$\vec B(x,y)$ is perpendicular to the $xy$-plane at any point in the plane}
\opt{B}{$|\vec B(x,y)|$ depends on $x$ and $y$ only through the radial distance $r=\sqrt{x^2+y^2}$}
\opt{C}{$|\vec B(x,y)|$ is non-zero at all points for $r<R$}
\opt{D}{$\vec B(x,y)$ points normally outward from the $xy$-plane for all the points between the two loops}
\end{qoptions}

\sect{SECTION-I(ii)}
\paraq{7}{8}
A special metal S conducts electricity without any resistance. A closed wire loop, made of S, does not allow any change in flux through itself by inducing a suitable current to generate a compensating flux. The induced current in the loop cannot decay due to its zero resistance. This current gives rise to a magnetic moment which in turn repels the source of magnetic field or flux. Consider such a loop, of radius $a$, with its center at the origin. A magnetic dipole of moment $m$ is brought along the axis of this loop from infinity to a point at distance $r$ ($\gg a$) from the center of the loop with its north pole always facing the loop, as shown in the figure below.\\
The magnitude of magnetic field of a dipole $m$, at a point on its axis at distance $r$, is $\dfrac{\mu_0}{2\pi}\dfrac{m}{r^3}$, where $\mu_0$ is the permeability of free space. The magnitude of the force between two magnetic dipoles with moments, $m_1$ and $m_2$, separated by a distance $r$ on the common axis, with their north poles facing each other, is $\dfrac{km_1m_2}{r^4}$, where $k$ is a constant of appropriate dimensions. The direction of this force is along the line joining the two dipoles.
\qfig{p2_phy_s2_para7_8_1}

\qn{7} When the dipole $m$ is placed at a distance $r$ from the center of the loop (as shown in the figure), the current induced in the loop will be proportional to
\begin{qoptions}
\opt{A}{$\dfrac{m}{r^3}$}
\opt{B}{$\dfrac{m^2}{r^2}$}
\opt{C}{$\dfrac{m}{r^2}$}
\opt{D}{$\dfrac{m^2}{r}$}
\end{qoptions}

\qn{8} The work done in bringing the dipole from infinity to a distance $r$ from the center of the loop by the given process is proportional to
\begin{qoptions}
\opt{A}{$\dfrac{m}{r^5}$}
\opt{B}{$\dfrac{m^2}{r^5}$}
\opt{C}{$\dfrac{m^2}{r^6}$}
\opt{D}{$\dfrac{m^2}{r^7}$}
\end{qoptions}

\paraq{9}{10}
A thermally insulating cylinder has a thermally insulating and frictionless movable partition in the middle, as shown in the figure below. On each side of the partition, there is one mole of an ideal gas, with specific heat at constant volume, $C_v=2R$. Here, $R$ is the gas constant. Initially, each side has a volume $V_0$ and temperature $T_0$. The left side has an electric heater, which is turned on at very low power to transfer heat $Q$ to the gas on the left side. As a result the partition moves slowly towards the right reducing the right side volume to $V_0/2$. Consequently, the gas temperatures on the left and the right sides become $T_L$ and $T_R$, respectively. Ignore the changes in the temperatures of the cylinder, heater and the partition.
\qfig{p2_phy_s2_para9_10_1}

\qn{9} The value of $\dfrac{T_R}{T_0}$ is
\begin{qoptions}
\opt{A}{$\sqrt2$}
\opt{B}{$\sqrt3$}
\opt{C}{2}
\opt{D}{3}
\end{qoptions}

\qn{10} The value of $\dfrac{Q}{RT_0}$ is
\begin{qoptions}
\opt{A}{$4\left(2\sqrt2+1\right)$}
\opt{B}{$4\left(2\sqrt2-1\right)$}
\opt{C}{$\left(5\sqrt2+1\right)$}
\opt{D}{$\left(5\sqrt2-1\right)$}
\end{qoptions}

\sect{SECTION-II(i)}
\paraq{1}{2}
A soft plastic bottle, filled with water of density 1 gm/cc, carries an inverted glass test-tube with some air (ideal gas) trapped as shown in the figure. The test-tube has a mass of 5 gm, and it is made of a thick glass of density 2.5 gm/cc. Initially the bottle is sealed at atmospheric pressure $p_0=10^5$ Pa so that the volume of the trapped air is $v_0=3.3$ cc. When the bottle is squeezed from outside at constant temperature, the pressure inside rises and the volume of the trapped air reduces. It is found that the test tube begins to sink at pressure $P_0+\Delta p$ without changing its orientation. At this pressure, the volume of the trapped air is $v_0-\Delta v$.\\
Let $\Delta v=X$ cc and $\Delta p=Y\times10^3$ Pa.
\qfig{p2_phy_s3_para1_2_1}

\qn{1} The value of $X$ is \underline{\hspace{2.5em}}.

\qn{2} The value of $Y$ is \underline{\hspace{2.5em}}.

\paraq{3}{4}
A pendulum consists of a bob of mass $m=0.1$ kg and a massless inextensible string of length $L=1.0$ m. It is suspended from a fixed point at height $H=0.9$ m above a frictionless horizontal floor. Initially, the bob of the pendulum is lying on the floor at rest vertically below the point of suspension. A horizontal impulse $P=0.2$ kg-m/s is imparted to the bob at some instant. After the bob slides for some distance, the string becomes taut and the bob lifts off the floor. The magnitude of the angular momentum of the pendulum about the point of suspension just before the bob lifts off is $J$ kg-m$^2$/s. The kinetic energy of the pendulum just after the lift-off is $K$ Joules.

\qn{3} The value of $J$ is \underline{\hspace{2.5em}}.

\qn{4} The value of $K$ is \underline{\hspace{2.5em}}.

\paraq{5}{6}
In a circuit, a metal filament lamp is connected in series with a capacitor of capacitance $C\ \mu$F across a 200 V, 50 Hz supply. The power consumed by the lamp is 500 W while the voltage drop across it is 100 V. Assume that there is no inductive load in the circuit. Take \textit{rms} values of the voltages. The magnitude of the phase-angle (in degrees) between the current and the supply voltage is $\varphi$. Assume, $\pi\sqrt3\approx5$.

\qn{5} The value of $C$ is \underline{\hspace{2.5em}}.

\qn{6} The value of $\varphi$ is \underline{\hspace{2.5em}}.

\sect{SECTION-II(ii)}
\qn{7} In order to measure the internal resistance $r_1$ of a cell of emf $E$, a meter bridge of wire resistance $R_0=50\ \Omega$, a resistance $R_0/2$, another cell of emf $E/2$ (internal resistance $r$) and a galvanometer $G$ are used in a circuit, as shown in the figure. If the null point is found at $\ell=72$ cm, then the value of $r_1=$ \underline{\hspace{2.5em}} $\Omega$.
\qfig{p2_phy_s4_q07_1}

\qn{8} The distance between two stars of masses $3M_S$ and $6M_S$ is $9R$. Here $R$ is the mean distance between the centers of the Earth and the Sun, and $M_S$ is the mass of the Sun. The two stars orbit around their common center of mass in circular orbits with period $nT$, where $T$ is the period of Earth's revolution around the Sun. The value of $n$ is \underline{\hspace{2.5em}}.

\qn{9} In a photoemission experiment, the maximum kinetic energies of photoelectrons from metals $P$, $Q$ and $R$ are $E_P$, $E_Q$ and $E_R$, respectively, and they are related by $E_P=2E_Q=2E_R$. In this experiment, the same source of monochromatic light is used for metals $P$ and $Q$ while a different source of monochromatic light is used for the metal $R$. The work functions for metals $P$, $Q$ and $R$ are 4.0 eV, 4.5 eV and 5.5 eV, respectively. The energy of the incident photon used for metal $R$, in eV, is \underline{\hspace{2.5em}}.

\raggedcolumns
\partsec{2}{CHEMISTRY}
\sect{SECTION-I(i)}
\qn{1} The reaction sequence(s) that would lead to $o$-xylene as the major product is (are)
\begin{qoptions}
\opt{A}{\qfigi{p2_chem_s1_q01_A}}
\opt{B}{\qfigi{p2_chem_s1_q01_B}}
\opt{C}{\qfigi{p2_chem_s1_q01_C}}
\opt{D}{\qfigi{p2_chem_s1_q01_D}}
\end{qoptions}

\qn{2} Correct option(s) for the following sequence of reactions is(are)
\qfig{p2_chem_s1_q02_1}
\begin{qoptions}
\opt{A}{$\mathbf Q=\mathrm{KNO_2}$, $\mathbf W=\mathrm{LiAlH_4}$}
\opt{B}{$\mathbf R=$ benzenamine, $\mathbf V=\mathrm{KCN}$}
\opt{C}{$\mathbf Q=\mathrm{AgNO_2}$, $\mathbf R=$ phenylmethanamine}
\opt{D}{$\mathbf W=\mathrm{LiAlH_4}$, $\mathbf V=\mathrm{AgCN}$}
\end{qoptions}

\qn{3} For the following reaction\\
$2\mathbf X+\mathbf Y\xrightarrow{\;k\;}\mathbf P$\\
the rate of reaction is $\dfrac{d[\mathrm P]}{dt}=k[\mathbf X]$. Two moles of $\mathbf X$ are mixed with one mole of $\mathbf Y$ to make 1.0 L of solution. At 50 s, 0.5 mole of $\mathbf Y$ is left in the reaction mixture. The correct statement(s) about the reaction is(are)\\
(Use: $\ln 2=0.693$)
\begin{qoptions}
\opt{A}{The rate constant, $k$, of the reaction is $13.86\times10^{-4}\ \mathrm{s^{-1}}$.}
\opt{B}{Half--life of $\mathbf X$ is 50 s.}
\opt{C}{At 50 s, $-\dfrac{d[\mathbf X]}{dt}=13.86\times10^{-3}\ \mathrm{mol\,L^{-1}\,s^{-1}}$.}
\opt{D}{At 100 s, $-\dfrac{d[\mathbf Y]}{dt}=3.46\times10^{-3}\ \mathrm{mol\,L^{-1}\,s^{-1}}$}
\end{qoptions}

\qn{4} Some standard electrode potentials at 298 K are given below:\\
\begin{tabular}{@{}ll@{}}
$\mathrm{Pb^{2+}/Pb}$ & $-0.13$ V\\
$\mathrm{Ni^{2+}/Ni}$ & $-0.24$ V\\
$\mathrm{Cd^{2+}/Cd}$ & $-0.40$ V\\
$\mathrm{Fe^{2+}/Fe}$ & $-0.44$ V
\end{tabular}\\[2pt]
To a solution containing 0.001 M of $\mathbf X^{2+}$ and 0.1 M of $\mathbf Y^{2+}$, the metal rods $\mathbf X$ and $\mathbf Y$ are inserted (at 298 K) and connected by a conducting wire. This resulted in dissolution of $\mathbf X$. The correct combination(s) of $\mathbf X$ and $\mathbf Y$, respectively, is (are)\\
(Given: Gas constant, $R=8.314\ \mathrm{J\,K^{-1}\,mol^{-1}}$, Faraday constant, $F=96500\ \mathrm{C\,mol^{-1}}$)
\begin{qoptions}
\opt{A}{Cd and Ni}
\opt{B}{Cd and Fe}
\opt{C}{Ni and Pb}
\opt{D}{Ni and Fe}
\end{qoptions}

\qn{5} The pair(s) of complexes wherein both exhibit tetrahedral geometry is(are)\\
(Note: py = pyridine\\
Given: Atomic numbers of Fe, Co, Ni and Cu are 26, 27, 28 and 29, respectively)
\begin{qoptions}
\opt{A}{$[\mathrm{FeCl_4}]^{-}$ and $[\mathrm{Fe(CO)_4}]^{2-}$}
\opt{B}{$[\mathrm{Co(CO)_4}]^{-}$ and $[\mathrm{CoCl_4}]^{2-}$}
\opt{C}{$[\mathrm{Ni(CO)_4}]$ and $[\mathrm{Ni(CN)_4}]^{2-}$}
\opt{D}{$[\mathrm{Cu(py)_4}]^{+}$ and $[\mathrm{Cu(CN)_4}]^{3-}$}
\end{qoptions}

\qn{6} The correct statement(s) related to oxoacids of phosphorous is(are)
\begin{qoptions}
\opt{A}{Upon heating, $\mathrm{H_3PO_3}$ undergoes disproportionation reaction to produce $\mathrm{H_3PO_4}$ and $\mathrm{PH_3}$.}
\opt{B}{While $\mathrm{H_3PO_3}$ can act as reducing agent, $\mathrm{H_3PO_4}$ cannot.}
\opt{C}{$\mathrm{H_3PO_3}$ is a monobasic acid.}
\opt{D}{The H atom of P--H bond in $\mathrm{H_3PO_3}$ is not ionizable in water.}
\end{qoptions}

\sect{SECTION-I(ii)}
\paraq{7}{8}
The amount of energy required to break a bond is same as the amount of energy released when the same bond is formed. In gaseous state, the energy required for \textit{homolytic cleavage} of a bond is called Bond Dissociation Energy (BDE) or Bond Strength. BDE is affected by $s$-character of the bond and the stability of the radicals formed. Shorter bonds are typically stronger bonds. BDEs for some bonds are given below :
\qfig{p2_chem_s2_para7_8_1}
$\mathrm{Cl-Cl(g)\to Cl^{\bullet}(g)+Cl^{\bullet}(g)}$\quad $\Delta H^\circ=58$ kcal mol$^{-1}$\\
$\mathrm{H_3C-Cl(g)\to H_3C^{\bullet}(g)+Cl^{\bullet}(g)}$\quad $\Delta H^\circ=85$ kcal mol$^{-1}$\\
$\mathrm{H-Cl(g)\to H^{\bullet}(g)+Cl^{\bullet}(g)}$\quad $\Delta H^\circ=103$ kcal mol$^{-1}$

\qn{7} Correct match of the \textbf{C}--\textbf{H} bonds (shown in bold) in Column \textbf{J} with their BDE in Column \textbf{K} is
\par\nobreak
\begin{center}\small
\begin{tabular}{|p{2.9cm}|p{2.9cm}|}
\hline
Column \textbf J\newline Molecule & Column \textbf K\newline BDE (kcal mol$^{-1}$)\\ \hline
(P) $\mathbf{H\text{--}C}\mathrm{H(CH_3)_2}$ & (i) 132\\ \hline
(Q) $\mathrm H\text{--}\mathbf{C}\mathrm H_2\mathrm{Ph}$ & (ii) 110\\ \hline
(R) $\mathbf{H\text{--}C}\mathrm{H{=}CH_2}$ & (iii) 95\\ \hline
(S) $\mathrm H\text{--}\mathbf C{\equiv}\mathrm{CH}$ & (iv) 88\\ \hline
\end{tabular}
\end{center}
\begin{qoptions}
\opt{A}{P -- iii, Q -- iv, R -- ii, S -- i}
\opt{B}{P -- i, Q -- ii, R -- iii, S -- iv}
\opt{C}{P -- iii, Q -- ii, R -- i, S -- iv}
\opt{D}{P -- ii, Q -- i, R -- iv, S -- iii}
\end{qoptions}

\qn{8} For the following reaction\\
$\mathrm{CH_4(g)+Cl_2(g)}\xrightarrow{\;\text{light}\;}\mathrm{CH_3Cl(g)+HCl(g)}$\\
the correct statement is
\begin{qoptions}
\opt{A}{Initiation step is exothermic with $\Delta H^\circ=-58$ kcal mol$^{-1}$}
\opt{B}{Propagation step involving $^{\bullet}\mathrm{CH_3}$ formation is exothermic with $\Delta H^\circ=-2$ kcal mol$^{-1}$.}
\opt{C}{Propagation step involving $\mathrm{CH_3Cl}$ formation is endothermic with $\Delta H^\circ=+27$ kcal mol$^{-1}$.}
\opt{D}{The reaction is exothermic with $\Delta H^\circ=-25$ kcal mol$^{-1}$.}
\end{qoptions}

\paraq{9}{10}
The reaction of $\mathrm{K_3[Fe(CN)_6]}$ with freshly prepared $\mathrm{FeSO_4}$ solution produces a dark blue precipitate called Turnbull's blue. Reaction of $\mathrm{K_4[Fe(CN)_6]}$ with the $\mathrm{FeSO_4}$ solution in complete absence of air produces a white precipitate \textbf{X}, which turns blue in air. Mixing the $\mathrm{FeSO_4}$ solution with $\mathrm{NaNO_3}$, followed by a slow addition of concentrated $\mathrm{H_2SO_4}$ through the side of the test tube produces a brown ring.

\qn{9} Precipitate \textbf{X} is
\begin{qoptions}
\opt{A}{$\mathrm{Fe_4[Fe(CN)_6]_3}$}
\opt{B}{$\mathrm{Fe[Fe(CN)_6]}$}
\opt{C}{$\mathrm{K_2Fe[Fe(CN)_6]}$}
\opt{D}{$\mathrm{KFe[Fe(CN)_6]}$}
\end{qoptions}

\qn{10} Among the following, the brown ring is due to the formation of
\begin{qoptions}
\opt{A}{$[\mathrm{Fe(NO)_2(SO_4)_2}]^{2-}$}
\opt{B}{$[\mathrm{Fe(NO)_2(H_2O)_4}]^{3+}$}
\opt{C}{$[\mathrm{Fe(NO)_4(SO_4)_2}]$}
\opt{D}{$[\mathrm{Fe(NO)(H_2O)_5}]^{2+}$}
\end{qoptions}

\sect{SECTION-II(i)}
\paraq{1}{2}
At 298 K, the limiting molar conductivity of a weak monobasic acid is $4\times10^2\ \mathrm{S\,cm^2\,mol^{-1}}$. At 298 K, for an aqueous solution of the acid the degree of dissociation of $\boldsymbol\alpha$ and the molar conductivity is $\mathbf y\times10^2\ \mathrm{S\,cm^2\,mol^{-1}}$. At 298 K, upon 20 times dilution with water, the molar conductivity of the solution becomes $3\mathbf y\times10^2\ \mathrm{S\,cm^2\,mol^{-1}}$.

\qn{1} The value of $\boldsymbol\alpha$ is \underline{\hspace{2.5em}}.

\qn{2} The value of $\mathbf y$ is \underline{\hspace{2.5em}}.

\paraq{3}{4}
Reaction of $\mathbf x$ g of Sn with HCl quantitatively produced a salt. Entire amount of the salt reacted with $\mathbf y$ g of nitrobenzene in the presence of required amount of HCl to produce 1.29 g of an organic salt (quantitatively).\\
(Use Molar masses (in g mol$^{-1}$) of H, C, N, O, Cl and Sn as 1, 12, 14, 16, 35 and 119, respectively).

\qn{3} The value of $\mathbf x$ is \underline{\hspace{2.5em}}.

\qn{4} The value of $\mathbf y$ is \underline{\hspace{2.5em}}.

\paraq{5}{6}
A sample (5.6 g) containing iron is completely dissolved in cold dilute HCl to prepare a 250 mL of solution. Titration of 25.0 mL of this solution requires 12.5 mL of 0.03 M $\mathrm{KMnO_4}$ solution to reach the end point. Number of moles of $\mathrm{Fe^{2+}}$ present in 250 mL solution is $\mathbf x\times10^{-2}$ (consider complete dissolution of $\mathrm{FeCl_2}$). The amount of iron present in the sample of $\mathbf{y\%}$ by weight.\\
(Assume : $\mathrm{KMnO_4}$ reacts only with $\mathrm{Fe^{2+}}$ in the solution Use : Molar mass of iron as 56 g mol$^{-1}$)

\qn{5} The value of $\mathbf x$ is \underline{\hspace{2.5em}}.

\qn{6} The value of $\mathbf y$ is \underline{\hspace{2.5em}}.

\sect{SECTION-II(ii)}
\qn{7} One mole of an ideal gas at 900 K, undergoes two reversible processes, \textbf{I} followed by \textbf{II}, as shown below. If the work done by the gas in the two processes are same, the value of $\ln\dfrac{V_3}{V_2}$ is \underline{\hspace{2.5em}}.
\qfig{p2_chem_s4_q07_1}
($U$ : internal energy, $S$ : entropy, $p$ : pressure, $V$ : volume, $R$ : gas constant)\\
(Given : molar heat capacity at constant volume, $C_{V,\mathrm m}$ of the gas is $\dfrac52R$)

\qn{8} Consider a helium (He) atom that absorbs a photon of wavelength 330 nm. The change in the velocity (in cm s$^{-1}$) of He atom after the photon absorption is \underline{\hspace{2.5em}}.\\
(Assume: Momentum is conserved when photon is absorbed.\\
Use: Planck constant $=6.6\times10^{-34}$ J s, Avogadro number $=6\times10^{23}\ \mathrm{mol^{-1}}$, Molar mass of He $=4\ \mathrm{g\,mol^{-1}}$)

\qn{9} Ozonolysis of $\mathrm{ClO_2}$ produces an oxide of chlorine. The average oxidation state of chlorine in this oxide is \underline{\hspace{2.5em}}.

\raggedcolumns
\partsec{3}{MATHEMATICS}
\sect{SECTION-I(i)}
\qn{1} Let\\
$S_1=\{(i,j,k):i,j,k\in\{1,2,\dots,10\}\}$,\\
$S_2=\{(i,j):1\le i<j+2\le10,$ $i,j\in\{1,2,\dots,10\}\}$,\\
$S_3=\{(i,j,k,\ell):1\le i<j<k<\ell,$ $i,j,k,\ell\in\{1,2,\dots,10\}\}$\\
and\\
$S_4=\{(i,j,k,\ell):i,j,k$ and $\ell$ are distinct elements in $\{1,2,\dots,10\}\}$.\\
If the total number of elements in the set $S_r$ is $n_r$, $r=1,2,3,4$, then which of the following statements is (are) \textbf{TRUE} ?
\begin{qoptions}
\opt{A}{$n_1=1000$}
\opt{B}{$n_2=44$}
\opt{C}{$n_3=220$}
\opt{D}{$\dfrac{n_4}{12}=420$}
\end{qoptions}

\qn{2} Consider a triangle $PQR$ having sides of lengths $p$, $q$ and $r$ opposite to the angles $P$, $Q$ and $R$, respectively. Then which of the following statements is (are) \textbf{TRUE} ?
\begin{qoptions}
\opt{A}{$\cos P\ge1-\dfrac{p^2}{2qr}$}
\opt{B}{$\cos R\ge\left(\dfrac{q-r}{p+q}\right)\cos P+\left(\dfrac{p-r}{p+q}\right)\cos Q$}
\opt{C}{$\dfrac{q+r}{p}<2\dfrac{\sqrt{\sin Q\sin R}}{\sin P}$}
\opt{D}{If $p<q$ and $p<r$, then $\cos Q>\dfrac pr$ and $\cos R>\dfrac pq$}
\end{qoptions}

\qn{3} Let $f:\left[-\dfrac\pi2,\dfrac\pi2\right]\to\mathbb R$ be a continuous function such that $f(0)=1$ and $\displaystyle\int_0^{\pi/3}f(t)\,dt=0$. Then which of the following statements is (are) \textbf{TRUE} ?
\begin{qoptions}
\opt{A}{The equation $f(x)-3\cos3x=0$ has at least one solution in $\left(0,\dfrac\pi3\right)$}
\opt{B}{The equation $f(x)-3\sin3x=-\dfrac6\pi$ has at least one solution in $\left(0,\dfrac\pi3\right)$}
\opt{C}{$\displaystyle\lim_{x\to0}\dfrac{x\int_0^xf(t)\,dt}{1-e^{x^2}}=-1$}
\opt{D}{$\displaystyle\lim_{x\to0}\dfrac{\sin x\int_0^xf(t)\,dt}{x^2}=-1$}
\end{qoptions}

\qn{4} For any real numbers $\alpha$ and $\beta$, let $y_{\alpha,\beta}(x)$, $x\in\mathbb R$, be the solution of the differential equation $\dfrac{dy}{dx}+\alpha y=xe^{\beta x}$, $y(1)=1$. Let $S=\{y_{\alpha,\beta}(x):\alpha,\beta\in\mathbb R\}$. Then which of the following functions belong(s) to the set $S$ ?
\begin{qoptions}
\opt{A}{$f(x)=\dfrac{x^2}{2}e^{-x}+\left(e-\dfrac12\right)e^{-x}$}
\opt{B}{$f(x)=-\dfrac{x^2}{2}e^{-x}+\left(e+\dfrac12\right)e^{-x}$}
\opt{C}{$f(x)=\dfrac{e^x}{2}\left(x-\dfrac12\right)+\left(e-\dfrac{e^2}{4}\right)e^{-x}$}
\opt{D}{$f(x)=\dfrac{e^x}{2}\left(\dfrac12-x\right)+\left(e+\dfrac{e^2}{4}\right)e^{-x}$}
\end{qoptions}

\qn{5} Let $O$ be the origin and $\overrightarrow{OA}=2\hat i+2\hat j+\hat k$, $\overrightarrow{OB}=\hat i-2\hat j+2\hat k$ and $\overrightarrow{OC}=\dfrac12\left(\overrightarrow{OB}-\lambda\overrightarrow{OA}\right)$ for some $\lambda>0$. If $\left|\overrightarrow{OB}\times\overrightarrow{OC}\right|=\dfrac92$, then which of the following statements is (are) \textbf{TRUE} ?
\begin{qoptions}
\opt{A}{Projection of $\overrightarrow{OC}$ on $\overrightarrow{OA}$ is $-\dfrac32$}
\opt{B}{Area of the triangle $OAB$ is $\dfrac92$}
\opt{C}{Area of the triangle $ABC$ is $\dfrac92$}
\opt{D}{The acute angle between the diagonals of the parallelogram with adjacent sides $\overrightarrow{OA}$ and $\overrightarrow{OC}$ is $\dfrac\pi3$}
\end{qoptions}

\qn{6} Let $E$ denote the parabola $y^2=8x$. Let $P=(-2,4)$, and let $Q$ and $Q'$ be two distinct points on $E$ such that the lines $PQ$ and $PQ'$ are tangents to $E$. Let $F$ be the focus of $E$. Then which of the following statements is (are) \textbf{TRUE} ?
\begin{qoptions}
\opt{A}{The triangle $PFQ$ is a right-angled triangle}
\opt{B}{The triangle $QPQ'$ is a right-angled triangle}
\opt{C}{The distance between $P$ and $F$ is $5\sqrt2$}
\opt{D}{$F$ lies on the line joining $Q$ and $Q'$}
\end{qoptions}

\sect{SECTION-I(ii)}
\paraq{7}{8}
Let $M=\{(x,y)\in\mathbb R\times\mathbb R:x^2+y^2\le r^2\}$, where $r>0$. Consider the geometric progression $a_n=\dfrac1{2^{n-1}}$, $n=1,2,3,\dots$. Let $S_0=0$ and, for $n\ge1$, let $S_n$ denote the sum of the first $n$ terms of this progression. For $n\ge1$, let $C_n$ denote the circle with center $(S_{n-1},0)$ and radius $a_n$, and $D_n$ denote the circle with center $(S_{n-1},S_{n-1})$ and radius $a_n$.

\qn{7} Consider $M$ with $r=\dfrac{1025}{513}$. Let $k$ be the number of all those circles $C_n$ that are inside $M$. Let $\ell$ be the maximum possible number of circles among these $k$ circles such that no two circles intersect. Then
\begin{qoptions}
\opt{A}{$k+2\ell=22$}
\opt{B}{$2k+\ell=26$}
\opt{C}{$2k+3\ell=34$}
\opt{D}{$3k+2\ell=40$}
\end{qoptions}

\qn{8} Consider $M$ with $r=\dfrac{(2^{199}-1)\sqrt2}{2^{198}}$. The number of all those circles $D_n$ that are inside $M$ is
\begin{qoptions}
\opt{A}{198}
\opt{B}{199}
\opt{C}{200}
\opt{D}{201}
\end{qoptions}

\paraq{9}{10}
Let $\psi_1:[0,\infty)\to\mathbb R$, $\psi_2:[0,\infty)\to\mathbb R$, $f:[0,\infty)\to\mathbb R$ and $g:[0,\infty)\to\mathbb R$ be functions such that\\
$f(0)=g(0)=0$,\\
$\psi_1(x)=e^{-x}+x$, $x\ge0$,\\
$\psi_2(x)=x^2-2x-2e^{-x}+2$, $x\ge0$,\\
$f(x)=\displaystyle\int_{-x}^{x}\left(|t|-t^2\right)e^{-t^2}\,dt$, $x>0$\\
and\\
$g(x)=\displaystyle\int_0^{x^2}\sqrt t\,e^{-t}\,dt$, $x>0$

\qn{9} Which of the following statements is \textbf{TRUE} ?
\begin{qoptions}
\opt{A}{$f\left(\sqrt{\ln3}\right)+g\left(\sqrt{\ln3}\right)=\dfrac13$}
\opt{B}{For every $x>1$, there exists an $\alpha\in(1,x)$ such that $\psi_1(x)=1+\alpha x$}
\opt{C}{For every $x>0$, there exists an $\beta\in(0,x)$ such that $\psi_2(x)=2x(\psi_1(\beta)-1)$}
\opt{D}{$f$ is an increasing function on the interval $\left[0,\dfrac32\right]$}
\end{qoptions}

\qn{10} Which of the following statements is \textbf{TRUE} ?
\begin{qoptions}
\opt{A}{$\psi_1(x)\le1$, for all $x>0$}
\opt{B}{$\psi_2(x)\le0$, for all $x>0$}
\opt{C}{$f(x)\ge1-e^{-x^2}-\dfrac23x^3+\dfrac25x^5$, for all $x\in\left(0,\dfrac12\right)$}
\opt{D}{$g(x)\le\dfrac23x^3-\dfrac25x^5+\dfrac17x^7$, for all $x\in\left(0,\dfrac12\right)$}
\end{qoptions}

\sect{SECTION-II(i)}
\paraq{1}{2}
Consider the region $R=\{(x,y)\in\mathbb R\times\mathbb R:x\ge0\text{ and }y^2\le4-x\}$. Let $\mathcal F$ be the family of all circles that are contained in $R$ and have centers on the $x$-axis. Let $C$ be the circle that has largest radius among the circles in $\mathcal F$. Let $(\alpha,\beta)$ be a point where the circle $C$ meets the curve $y^2=4-x$.

\qn{1} The radius of the circle $C$ is \underline{\hspace{2.5em}}.

\qn{2} The value of $\alpha$ is \underline{\hspace{2.5em}}.

\paraq{3}{4}
Let $f_1:(0,\infty)\to\mathbb R$ and $f_2:(0,\infty)\to\mathbb R$ be defined by $f_1(x)=\displaystyle\int_0^x\prod_{j=1}^{21}(t-j)^j\,dt$, $x>0$ and $f_2(x)=98(x-1)^{50}-600(x-1)^{49}+2450$, $x>0$, where, for any positive integer $n$ and real numbers $a_1,a_2,\dots,a_n$, $\displaystyle\prod_{i=1}^{n}a_i$ denotes the product of $a_1,a_2,\dots,a_n$. Let $m_i$ and $n_i$, respectively, denote the number of points of local minima and the number of points of local maxima of function $f_i$, $i=1,2$, in the interval $(0,\infty)$.

\qn{3} The value of $2m_1+3n_1+m_1n_1$ is \underline{\hspace{2.5em}}.

\qn{4} The value of $6m_2+4n_2+8m_2n_2$ is \underline{\hspace{2.5em}}.

\paraq{5}{6}
Let $g_i:\left[\dfrac\pi8,\dfrac{3\pi}8\right]\to\mathbb R$, $i=1,2$, and $f:\left[\dfrac\pi8,\dfrac{3\pi}8\right]\to\mathbb R$ be functions such that $g_1(x)=1$, $g_2(x)=|4x-\pi|$ and $f(x)=\sin^2x$, for all $x\in\left[\dfrac\pi8,\dfrac{3\pi}8\right]$.\\
Define $S_i=\displaystyle\int_{\pi/8}^{3\pi/8}f(x)\cdot g_i(x)\,dx$, $i=1,2$

\qn{5} The value of $\dfrac{16S_1}{\pi}$ is \underline{\hspace{2.5em}}.

\qn{6} The value of $\dfrac{48S_2}{\pi^2}$ is \underline{\hspace{2.5em}}.

\sect{SECTION-II(ii)}
\qn{7} A number is chosen at random from the set $\{1,2,3,\dots,2000\}$. Let $p$ be the probability that the chosen number is a multiple of 3 or a multiple of 7. Then the value of $500p$ is \underline{\hspace{2.5em}}.

\qn{8} Let $E$ be the ellipse $\dfrac{x^2}{16}+\dfrac{y^2}{9}=1$. For any three distinct points $P$, $Q$ and $Q'$ on $E$, let $M(P,Q)$ be the mid-point of the line segment joining $P$ and $Q$, and $M(P,Q')$ be the mid-point of the line segment joining $P$ and $Q'$. Then the maximum possible value of the distance between $M(P,Q)$ and $M(P,Q')$, as $P$, $Q$ and $Q'$ vary on $E$, is \underline{\hspace{2.5em}}.

\qn{9} For any real number $x$, let $[x]$ denote the largest integer less than or equal to $x$. If $I=\displaystyle\int_0^{10}\left[\sqrt{\dfrac{10x}{x+1}}\right]dx$, then the value of $9I$ is \underline{\hspace{2.5em}}.

\end{multicols}
\end{document}
